%======================================================================
\section{Log-signature compound BSDE with stop-gradient interfaces}\label{sec:method}
%======================================================================

The proposed method comprises four components: a two-grid simulation with local log-signatures, a shared Neural RDE encoder, layer-specific heads with a forward Euler recursion, and a training loss with detached interface targets. \Cref{fig:framework} illustrates how these components are organized, and \cref{alg:train} summarizes the training loop.

\begin{figure}[htbp]
\centering
\resizebox{\textwidth}{!}{% Framework figure: (a) time grids, (b) shared encoder and layer heads,
% (c) stop-gradient compound loss.  Simplified from the proposal figures.
\begin{tikzpicture}[
  font=\footnotesize,
  >={Stealth[length=2mm]},
  plab/.style={font=\small\bfseries, text=cInk, anchor=west, inner sep=0pt},
  rowlab/.style={font=\scriptsize, text=cInk, anchor=east, inner sep=2pt},
  note/.style={font=\scriptsize, text=cInk!65},
  ebox/.style={draw=cInk!60, fill=cStop!10, rounded corners=2pt, align=center,
               inner sep=3pt, minimum height=6.5mm, text width=2.7cm, font=\scriptsize},
  ar/.style={->, draw=cInk!65, line width=0.6pt},
  hd/.style={->, draw=cStop!80!black, line width=0.8pt},
  lay/.style={draw=cInk!80, line width=1pt},
  dotn/.style={circle, fill=cInk, inner sep=1.3pt},
  rbox/.style={draw=cInk!55, fill=white, rounded corners=2pt, align=center,
               inner sep=3pt, font=\scriptsize},
  lbox/.style={draw=cInk!70, fill=cStop!10, rounded corners=2pt, align=center,
               inner sep=4pt, font=\scriptsize},
  grad/.style={->, draw=cJoint!85!black, line width=0.8pt, dash pattern=on 3pt off 1.5pt},
  gstop/.style={draw=cJoint!85!black, line width=0.8pt, dash pattern=on 3pt off 1.5pt,
                -{Bar[width=3mm,line width=1.4pt]}},
]

%======================================================================
% (a) Time grids
%======================================================================
\begin{scope}[shift={(0,0)}, x=0.31cm]
  \node[plab] at (-8.3,2.75) {(a) Time grids};
  \fill[cStop!12] (4,-0.12) rectangle (8,1.05);
  % fine grid
  \draw[->, cInk!70] (-0.5,0) -- (33,0);
  \foreach \n in {0,...,32}{\draw[cInk!45] (\n,-0.07) -- (\n,0.07);}
  \node[rowlab] at (-0.6,0) {fine grid $t_n$};
  % learning grid
  \draw[->, cInk!70] (-0.5,0.9) -- (33,0.9);
  \foreach \n in {0,4,...,32}{\draw[thick, cStop!80!black] (\n,0.77) -- (\n,1.03);}
  \foreach \n/\l in {0/0,4/1,8/2,12/3,16/4,20/5,24/6,28/7,32/8}{
    \node[font=\tiny, anchor=north, inner sep=1.5pt] at (\n,0.76) {$u_{\l}$};}
  \node[rowlab] at (-0.6,0.9) {learning grid $u_i$};
  % exercise dates
  \draw[->, cInk!70] (-0.5,1.8) -- (33,1.8);
  \foreach \n in {0,12,20,32}{\fill[cInk] (\n,1.8) circle (1.8pt);}
  \foreach \n/\l in {0/{T_0=0},12/{T_1},20/{T_2},32/{T_3=T}}{
    \node[font=\tiny, anchor=south, inner sep=2pt] at (\n,1.84) {$\l$};}
  \node[rowlab] at (-0.6,1.8) {exercise dates $T_j$};
  \foreach \n in {0,4,...,32}{\draw[densely dashed, cInk!20] (\n,0.08) -- (\n,0.76);}
  \foreach \n in {0,12,20,32}{\draw[densely dashed, cInk!20] (\n,1.04) -- (\n,1.76);}
  % one learning interval
  \draw[->, cStop!80!black, thick] (6,-0.45) -- (6,0.7);
  \node[note, anchor=west, text=cStop!70!black] at (6.3,-0.38)
    {$\ell_1=\LogSig^M(X)_{u_1,u_2}$, \quad $\Delta W_1=W_{u_2}-W_{u_1}$};
\end{scope}

%======================================================================
% (b) Shared encoder and layer heads
%======================================================================
\begin{scope}[shift={(0,-1.6)}]
  \node[plab] at (-2.57,0.75) {(b) Shared encoder and layer heads};
  \node[ebox] (E1) at (-0.6,-0.55) {fine-grid path $\Ipl S^{\mathrm f}$};
  \node[ebox] (E2) at (-0.6,-1.45) {local log-signatures $\ell_i$};
  \node[ebox] (E3) at (-0.6,-2.35) {shared NRDE \eqref{eq:logode}};
  \node[ebox] (E4) at (-0.6,-3.25) {hidden states $H_{u_i}$};
  \draw[ar] (E1) -- (E2); \draw[ar] (E2) -- (E3); \draw[ar] (E3) -- (E4);
  % heads bus
  \draw[cStop!80!black, line width=0.8pt] (1.55,-0.85) -- (1.55,-3.25);
  \draw[hd] (E4.east) -- (1.55,-3.25);
  \node[note, anchor=north, align=center] at (1.55,-3.4) {head input\\ $\xi_i=(H_{u_i},X_{u_i})$};
  % layers
  \foreach \y/\lab/\hl/\ta/\tb/\val in {%
     -0.85/{layer 1}/{y_0,\ \psi^Z_1}/{T_0}/{T_1}/{\hY_{1,T_1}},
     -1.85/{layer 2}/{\psi^Y_2,\ \psi^Z_2}/{T_1}/{T_2}/{\hY_{2,T_2}},
     -3.25/{layer $K$}/{\psi^Y_K,\ \psi^Z_K}/{T_{K-1}}/{T_K}/{\hY_{K,T}}}{
    \draw[hd] (1.55,\y) -- (3.6,\y);
    \node[font=\scriptsize, anchor=south, inner sep=1.5pt] at (2.55,\y) {$\hl$};
    \draw[lay] (3.65,\y) -- (9.2,\y);
    \foreach \x in {3.65,4.45,...,9.2}{\draw[cInk!50] (\x,\y-0.07) -- (\x,\y+0.07);}
    \node[dotn] at (3.65,\y) {}; \node[dotn] at (9.2,\y) {};
    \node[note, anchor=south west, inner sep=2pt] at (3.7,\y+0.03) {\lab};
    \node[font=\tiny, anchor=north, inner sep=2pt] at (3.65,\y-0.08) {$\ta$};
    \node[font=\tiny, anchor=north, inner sep=2pt] at (9.2,\y-0.08) {$\tb$};
    \node[font=\scriptsize, anchor=west, inner sep=3pt] at (9.3,\y) {$\val$};}
  \node at (6.4,-2.5) {$\vdots$};
  \draw[->, cInk!65, line width=0.8pt] (3.65,-0.2) -- (9.2,-0.2);
  \node[font=\scriptsize, anchor=south, inner sep=1.5pt] at (6.42,-0.2)
    {forward recursion \eqref{eq:scheme}: $\hY_{j,u_{i+1}}=\hY_{j,u_i}+\hZ_{u_i}^\top\Delta W_i$};
\end{scope}

%======================================================================
% (c) Stop-gradient compound loss
%======================================================================
\begin{scope}[shift={(0,-6.5)}]
  \node[plab] at (-2.57,0.25) {(c) Stop-gradient compound loss};
  % staircase of layers
  \draw[lay, ->] (-1.2,-0.3) -- (2.6,-0.3);
  \draw[lay, ->] (2.6,-1.0) -- (6.4,-1.0);
  \draw[lay, ->] (6.4,-1.7) -- (10.2,-1.7);
  \foreach \p in {(-1.2,-0.3),(2.6,-0.3),(2.6,-1.0),(6.4,-1.0),(6.4,-1.7),(10.2,-1.7)}{\node[dotn] at \p {};}
  \node[font=\scriptsize, anchor=east] at (-1.25,-0.3) {$y_0$};
  \node[font=\scriptsize, anchor=south east, inner sep=1pt] at (2.55,-0.25) {$\hY_{1,T_1}$};
  \node[font=\scriptsize, anchor=south east, inner sep=1pt] at (6.35,-0.95) {$\hY_{2,T_2}$};
  \node[font=\scriptsize, anchor=south, inner sep=2pt] at (10.2,-1.65) {$\hY_{K,T}$};
  \node[font=\scriptsize, anchor=north east, inner sep=1pt] at (2.55,-1.05) {$\hY_{2,T_1}$};
  \node[font=\scriptsize, anchor=north east, inner sep=1pt] at (6.35,-1.75) {$\hY_{3,T_2}$};
  % residual boxes
  \node[rbox] (e1) at (2.6,-2.75) {$e^{\sg}_1=\hY_{1,T_1}-\max\{\Phi^{\mathrm f}_1,\sg[\hY_{2,T_1}]\}$};
  \node[rbox] (e2) at (7.4,-2.75) {$e^{\sg}_2=\hY_{2,T_2}-\max\{\Phi^{\mathrm f}_2,\sg[\hY_{3,T_2}]\}$};
  \node[rbox] (eK) at (11.6,-2.75) {$e_K=\hY_{K,T}-\Phi^{\mathrm f}_K$};
  % layer 1 end value: forward and gradient
  \draw[ar] (2.5,-0.36) to[bend right=40] ($(e1.north)+(-1.6,0)$);
  \draw[grad] ($(e1.north)+(-1.9,0)$) to[bend left=40] (2.42,-0.3);
  \node[font=\tiny, text=cJoint!85!black, anchor=east, align=right] at (0.4,-1.2) {gradient\\ updates layer 1};
  % continuation value: forward, gradient blocked
  \draw[ar] (2.6,-1.06) -- ($(e1.north)+(0,0)$);
  \draw[gstop] ($(e1.north)+(0.35,0)$) -- ++(0,0.75);
  \node[font=\tiny, text=cJoint!85!black, anchor=west, align=left] at ($(e1.north)+(0.45,0.45)$) {gradient\\ blocked};
  \draw[ar] (6.4,-1.06) -- (6.4,-2.45);
  \draw[ar] (10.2,-1.76) -- ($(eK.north)+(-1.4,0)$);
  % loss
  \node[lbox] (L) at (5.0,-4.0) {$\cL_\omega(\theta)=\sum_{j=1}^K\omega_j\E\abs{e_j}^2$; \ one Adam step updates $y_0$, $\psi^Y_j$, $\psi^Z_j$ and the encoder};
  \draw[ar] (e1.south) -- (e1.south |- L.north);
  \draw[ar] (e2.south) -- (e2.south |- L.north);
  \draw[ar] (eK.south) |- (L.east);
\end{scope}
\end{tikzpicture}
}
\caption{Overview of the method: (a) time grids, (b) shared encoder and layer heads, (c) stop-gradient compound loss.}
\label{fig:framework}
\end{figure}

\subsection{Grids and local log-signatures}\label{sec:grids}

We use two nested time grids: a fine grid $t_n=n\hf$, on which asset paths are simulated and payoffs are monitored, and a coarser learning grid $u_i=ih$ with $h=T/N$, on which the encoder and the BSDE recursion operate. The learning grid contains every exercise date, and the refinement ratio $m=h/\hf$ is an integer. The simulated prices $S^{\mathrm f}_{t_n}$ are driven by the increments of $W$ on the fine grid, exactly for geometric Brownian motion and by the Euler scheme otherwise, and $\Ipl S^{\mathrm f}$ denotes their piecewise-linear interpolation. Accordingly, the monitored payoffs and links are given by
\[
\Phi^{\mathrm f}_j=\Phi_j\left(\Ipl S^{\mathrm f}|_{[0,T_j]}\right),\quad g^{\mathrm f}_j(y)=g_j\left(\Ipl S^{\mathrm f}|_{[0,T_j]},y\right).
\]

The encoder takes as input the normalized time-augmented path $X_t=\left(t/T,\,D^{-1}(\Ipl S^{\mathrm f}_t-\mu)\right)\in\R^{1+d}$. The constants $\mu\in\R^d$ and the diagonal matrix $D$ are estimated once on a pilot sample and then frozen. The signature of $X$ on $[s,t]$ truncated at level $M$ collects the iterated integrals
\[
\Sig^{M}(X)_{s,t}=\left(1,\ \int_{s<v<t}\dd X_{v},\ \int_{s<v_1<v_2<t}\dd X_{v_1}\otimes\dd X_{v_2},\ \dots\right)
\]
up to order $M$. Chen's identity $\Sig^M(X)_{0,t}=\Sig^M(X)_{0,s}\otimes\Sig^M(X)_{s,t}$ assembles the signature of the whole path from those of its pieces \cite{lyons1998}. The log-signature $\LogSig^M=\log_\otimes\Sig^M$ encodes the same information in $q_M$ coordinates, far fewer than the signature. The encoder uses the local log-signatures
\begin{equation}\label{eq:local-logsig}
\ell_i=\LogSig^M(X)_{u_i,u_{i+1}}\in\R^{q_M},\quad i=0,\dots,N-1 .
\end{equation}
The first level records the increments of time and price, while the second level captures the signed L\'evy areas between coordinates. Since the L\'evy area between price coordinate $a$ and normalized time differs from $T^{-1}\int_{u_i}^{u_{i+1}}(X^a_t-X^a_{u_i})\dd t$ only by a multiple of the price increment, the level-two log-signature already encodes the local running integral of the price; \cref{lem:information} makes this precise.

\subsection{Shared Neural RDE encoder}\label{sec:encoder}

The hidden state $H_t\in\R^p$ starts from $H_0=\rho_\theta(X_0)$ and follows the log-ODE
\begin{equation}\label{eq:logode}
\frac{\dd H_s}{\dd s}=G_\theta(H_s)\,\frac{\ell_i}{h},\quad s\in[u_i,u_{i+1}],
\end{equation}
where $G_\theta:\R^p\to\R^{p\times q_M}$ is a neural vector field \cite{morrill2021}. A fixed number of Runge--Kutta sub-steps defines the update map $H_{u_{i+1}}=R_\theta(H_{u_i},\ell_i)$. A single state sequence $(H_{u_i})_{i\le N}$ extends across all layers without being reset at the exercise dates, and the encoder parameters are shared by every layer.

\subsection{Heads and forward recursion}\label{sec:heads}

Let $\xi_i=(H_{u_i},X_{u_i})$ denote the input of the heads at $u_i$, which contains the hidden state, the normalized price, and the time. The price at time zero is a trainable scalar $y_0$. For $j\ge2$, a value head outputs the initial value of layer $j$, and on every layer a hedge head outputs the integrand:
\[
\hY_{j,T_{j-1}}=\psi^Y_{j}(\xi_{i}),\quad u_i=T_{j-1},\quad
\hZ_{u_i}=\psi^Z_{j}(\xi_i),\quad u_i\in[T_{j-1},T_j).
\]
Within each layer, the value follows the forward Euler recursion
\begin{equation}\label{eq:scheme}
\hY_{1,0}=y_0,\quad \hY_{j,u_{i+1}}=\hY_{j,u_i}+\hZ_{u_i}^\top\Delta W_i,\quad u_i\in[T_{j-1},T_j),
\end{equation}
where $\Delta W_i=W_{u_{i+1}}-W_{u_i}$. Between grid points, we set $\hY_{j,t}=\hY_{j,T_{j-1}}+\int_{T_{j-1}}^t\hZ_s^\top\dd W_s$ with $\hZ$ piecewise constant. In the implementation, undiscounted values follow the step $(1+rh)\hY_{j,u_i}+\hZ_{u_i}^\top\Delta W_i$, and \cref{lem:discount} shows that the discrete discount factor $(1+rh)^{-i}$ maps this step exactly onto \eqref{eq:scheme}. The interface and terminal residuals are defined by
\begin{equation}\label{eq:residual}
e_j=\hY_{j,T_j}-g^{\mathrm f}_j\left(\hY_{j+1,T_j}\right),\quad j<K,\quad e_K=\hY_{K,T}-\Phi^{\mathrm f}_K .
\end{equation}

By construction, the network is non-anticipative. To make this precise, let $\cG^{\pi}_{u_i}=\sigma(S^{\mathrm f}_{t_n}:t_n\le u_i)$ denote the fine-grid filtration and $\cG^{M}_{u_i}=\sigma(\ell_0,\dots,\ell_{i-1})$ the encoder filtration.

\begin{lemma}[non-anticipation]\label{lem:nonanticipation}
For every parameter $\theta$, every level $M$, and every number of Runge--Kutta sub-steps:
\begin{enumerate}[label=(\roman*)]
\item $H_{u_i}$ is $\cG^{M}_{u_i}$-measurable and $\cG^{M}_{u_i}\subseteq\cG^{\pi}_{u_i}\subseteq\cF_{u_i}$;
\item $\hZ_{u_i}$ is $\cG^{M}_{u_i}$-measurable and $\hY_{j,T_{j-1}}$ is $\cG^{M}_{T_{j-1}}$-measurable;
\item if $\E\abs{\hY_{j,T_{j-1}}}^2+h\sum_{u_i\in[T_{j-1},T_j)}\E\abs{\hZ_{u_i}}^2<\infty$, then $t\mapsto\hY_{j,t}$ is a square-integrable martingale on $[T_{j-1},T_j]$.
\end{enumerate}
\end{lemma}

\begin{proof}[Proof sketch]
Each $\ell_k$ is a Borel function of the fine-grid prices on $[u_k,u_{k+1}]$, and $H_{u_{i+1}}$ is a Borel function of $(H_{u_i},\ell_i)$, so induction over $i$ yields (i). Part (ii) follows because the normalized price $X_{u_i}$ is determined by the level-one increments $\ell_0,\dots,\ell_{i-1}$. Part (iii) is a consequence of $\E[\Delta W_i\mid\cF_{u_i}]=0$ and It\^o's isometry. See \cref{app:nonanticipation}.
\end{proof}

\begin{remark}\label{rem:nonanticipation}
The network value is a martingale on every layer for every choice of parameters, trained or not, and this is the only property of the architecture used in \cref{sec:analysis}. The lemma dictates two design rules: the hedge head at $u_i$ must not access the current increment $\ell_i$, and the normalization statistics $(\mu,D)$ must not be computed on the active training batch. Violating either rule induces a look-ahead bias with $\E[\hZ_{u_i}^\top\Delta W_i]\ne0$, which invalidates both the a posteriori estimate and the dual upper bound.
\end{remark}

\subsection{Stop-gradient interface targets}\label{sec:loss}

The training objective is the weighted endpoint loss
\begin{equation}\label{eq:loss}
\cL_\omega(\theta)=\sum_{j=1}^K\omega_j\,\E\abs{e_j}^2,\quad \omega_j>0,
\end{equation}
which is estimated on fresh batches of simulated paths. We write $\cL(\theta)$ for the unweighted loss with $\omega_j=1$; the experiments use $\omega_j=1/K$. The residual $e_j$ depends on the parameters through both sides of the interface: the left side $\hY_{j,T_j}$ depends on the heads of layer $j$ and on the encoder, whereas the target $g^{\mathrm f}_j(\hY_{j+1,T_j})$ depends on the value head of layer $j+1$. \emph{Joint training} \cite{huang2026} differentiates both sides.

\emph{Stop-gradient training} detaches the target instead. Denoting by $\sg[\cdot]$ the stop-gradient operator, which acts as the identity in the forward pass but has zero derivative in the backward pass, we define the detached residual
\begin{equation}\label{eq:sg-residual}
e^{\sg}_j=\hY_{j,T_j}-g^{\mathrm f}_j\left(\sg[\hY_{j+1,T_j}]\right).
\end{equation}
Write $\cL^{\sg}_\omega$ for \eqref{eq:loss} with $e^{\sg}_j$ in place of $e_j$. The case $K=2$ with the Bermudan link shows the difference. Let $\theta^Y_2$ denote the parameters of the value head $\psi^Y_2$, and write $v=\hY_{2,T_1}$ for its output. The value $v$ enters both residuals: as the initial value of layer $2$ in $e_2=v+\sum_{u_i\in[T_1,T)}\hZ_{u_i}^\top\Delta W_i-\Phi^{\mathrm f}_2$, and as the continuation value in $e_1=\hY_{1,T_1}-\max\{\Phi^{\mathrm f}_1,v\}$. By the chain rule,
\begin{equation}\label{eq:grad-K2}
\begin{aligned}
\nabla_{\theta^Y_2}\cL_\omega&=2\,\E\Bigl[\bigl(\omega_2e_2-\omega_1e_1\1\{v>\Phi^{\mathrm f}_1\}\bigr)\nabla_{\theta^Y_2}v\Bigr] &&\text{(joint)},\\
\nabla_{\theta^Y_2}\cL^{\sg}_\omega&=2\,\E\bigl[\omega_2e_2\,\nabla_{\theta^Y_2}v\bigr] &&\text{(stop-gradient)}.
\end{aligned}
\end{equation}
So the value head $\psi^Y_2$ is updated using only the second-layer residual $e_2$. Now suppose that the value head permits local perturbations $v+\varepsilon h$ for every bounded $\cG^M_{T_1}$-measurable function $h$, changing the residual to $e_2+\varepsilon h$. At a stationary point, $0=\left.\frac{\dd}{\dd\varepsilon}\omega_2\E[(e_2+\varepsilon h)^2]\right|_{\varepsilon=0}=2\omega_2\E[e_2h]$. Because this equality holds for every such $h$, it implies $\E[e_2\mid\cG^M_{T_1}]=0$. Hence, by the martingale property of layer $2$ (\cref{lem:nonanticipation}(iii)), we have $v=\E[\Phi^{\mathrm f}_2\mid\cG^M_{T_1}]$. However, under joint training, we have
\[
v-\E[\Phi^{\mathrm f}_2\mid\cG^M_{T_1}]=\frac{\omega_1}{\omega_2}\E\left[e_1\1\{v>\Phi^{\mathrm f}_1\}\bigm|\cG^M_{T_1}\right].
\]
Thus, the first-layer residual can shift $v$ away from the
conditional expectation of the future payoff. On paths where
$v>\Phi^{\mathrm f}_1$, the residual is $e_1=\hY_{1,T_1}-v$. If $e_1>0$, the first-layer prediction exceeds $v$, so increasing $v$ reduces the first-layer squared residual locally. This will accumulate errors through layers. Stop-gradient training removes this contribution to the update of the second-layer value head.  \cref{thm:stationary} states this observation detailedly for general $K$.

This construction transfers the semi-gradient rule of temporal-difference learning and the target network of deep $Q$-learning \cite{sutton2018,mnih2015} to the compound BSDE. Two related variants appear in \cref{sec:experiments}. The exponential moving average (EMA) variant evaluates the target with a slowly updated moving average of the parameters, and layerwise training uses the same detached target but trains the layers one at a time, backward from $T$. \Cref{fig:interface} compares the gradient paths.

\begin{figure}[htbp]
\centering
\resizebox{\textwidth}{!}{% Interface training variants.  Top: the forward computation shared by all
% variants.  Bottom: where the gradient of |e_j|^2 goes in each variant.
\begin{tikzpicture}[
    font=\footnotesize,
    >={Stealth[length=2mm]},
    param/.style={draw, rounded corners=2pt, minimum width=9mm, minimum height=6mm, inner sep=2pt},
    pj/.style={param, fill=cStop!14},
    pn/.style={param, fill=cFull!22},
    pema/.style={param, fill=cFull!10, dashed},
    pfrozen/.style={param, fill=cInk!12},
    res/.style={draw, circle, fill=white, inner sep=1pt, minimum size=7mm},
    tgt/.style={draw, rounded corners=6pt, fill=white, inner sep=2.5pt},
    f/.style={->, cInk, line width=0.6pt},
    g/.style={->, cJoint!85!black, line width=1pt, dash pattern=on 3pt off 1.5pt},
    gstop/.style={cJoint!85!black, line width=1pt, dash pattern=on 3pt off 1.5pt,
                  -{Bar[width=3mm,line width=1.4pt]}},
    ttl/.style={font=\footnotesize\bfseries, anchor=south, inner sep=1pt},
    nt/.style={font=\scriptsize, text=cInk!75, anchor=north, align=center, inner sep=1pt},
  ]

% ---------------- shared forward computation ----------------
\node[pj]  (thj) at (0,0)    {$\theta_j$};
\node      (yj)  at (2.2,0)  {$\hY_{j,T_j}$};
\node[res] (e)   at (4.3,0)  {$e_j$};
\node[tgt] (mx)  at (6.6,0)  {$\max\{\Phi^{\mathrm f}_j,\,\cdot\,\}$};
\node      (yn)  at (9.1,0)  {$\hY_{j+1,T_j}$};
\node[pn]  (thn) at (11.3,0) {$\theta_{j+1}$};
\draw[f] (thj) -- (yj); \draw[f] (yj) -- (e); \draw[f] (mx) -- (e);
\draw[f] (yn) -- (mx); \draw[f] (thn) -- (yn);
\node[font=\scriptsize, text=cInk!75, anchor=south] at (5.65,0.38)
  {shared forward computation (all variants)};
\draw[f] (12.3,0.75) -- ++(0.6,0) node[right, font=\scriptsize, text=cInk] {forward};
\draw[g] (12.3,0.35) -- ++(0.6,0) node[right, font=\scriptsize, text=cInk] {gradient};
\draw[gstop] (12.3,-0.05) -- ++(0.6,0) node[right, font=\scriptsize, text=cInk, xshift=1mm] {blocked};

% ---------------- the four gradient patterns ----------------
\foreach \x/\ttl in {0/{(a) Joint \cite{huang2026}}, 3.75/{(b) Stop-gradient}, 7.5/{(c) EMA target}, 11.25/{(d) Layerwise}}{
  \node[ttl] at (\x+1.25,-1.45) {\ttl};
}
% (a) joint
\begin{scope}[shift={(0,-2.2)}]
  \node[pj] (a1) at (0,0) {$\theta_j$}; \node[res] (a2) at (1.25,0) {$e_j$}; \node[pn] (a3) at (2.5,0) {$\theta_{j+1}$};
  \draw[g] (a2) -- (a1); \draw[g] (a2) -- (a3);
  \node[nt] at (1.25,-0.5) {target is differentiated;\\ all layers in one step};
\end{scope}
% (b) stop-gradient
\begin{scope}[shift={(3.75,-2.2)}]
  \node[pj] (b1) at (0,0) {$\theta_j$}; \node[res] (b2) at (1.25,0) {$e_j$}; \node[pn] (b3) at (2.5,0) {$\theta_{j+1}$};
  \draw[g] (b2) -- (b1); \draw[gstop] (b2) -- ($(b2)!0.6!(b3)$);
  \node[nt] at (1.25,-0.5) {target $\sg[\hY_{j+1,T_j}]$;\\ all layers in one step};
\end{scope}
% (c) EMA target
\begin{scope}[shift={(7.5,-2.2)}]
  \node[pj] (c1) at (0,0) {$\theta_j$}; \node[res] (c2) at (1.25,0) {$e_j$}; \node[pema] (c3) at (2.5,0) {$\bar\theta_{j+1}$};
  \draw[g] (c2) -- (c1); \draw[gstop] (c2) -- ($(c2)!0.6!(c3)$);
  \node[nt] at (1.25,-0.5) {target from lagged copy\\ $\bar\theta\leftarrow\beta\bar\theta+(1-\beta)\theta$};
\end{scope}
% (d) layerwise
\begin{scope}[shift={(11.25,-2.2)}]
  \node[pj] (d1) at (0,0) {$\theta_j$}; \node[res] (d2) at (1.25,0) {$e_j$}; \node[pfrozen] (d3) at (2.5,0) {$\theta^\star_{j+1}$};
  \draw[g] (d2) -- (d1); \draw[gstop] (d2) -- ($(d2)!0.6!(d3)$);
  \node[nt] at (1.25,-0.5) {$\theta^\star_{j+1}$ trained first, frozen;\\ one layer per stage};
\end{scope}
\end{tikzpicture}
}
\caption{Interface training variants. Top: forward computation. Bottom: gradient of $|e_j|^2$.}
\label{fig:interface}
\end{figure}

\begin{algorithm}[htbp]
\caption{Stop-gradient training of the log-signature compound BSDE}
\label{alg:train}
\begin{algorithmic}[1]
\Require grids $(t_n)$, $(u_i)\ni T_j$; level $M$; weights $\omega_j$; batch size $B$; number of updates $n_{\mathrm{it}}$
\For{$n=1,\dots,n_{\mathrm{it}}$}
  \State simulate $B$ fine-grid paths $S^{\mathrm f}$ with increments $\Delta W$
  \State compute $\ell_i$ by \eqref{eq:local-logsig} and $H_{u_{i+1}}=R_\theta(H_{u_i},\ell_i)$ for $i<N$
  \For{$j=1,\dots,K$}
    \State $\hY_{j,T_{j-1}}\gets y_0$ if $j=1$, else $\psi^Y_j(\xi_i)$ with $u_i=T_{j-1}$
    \State run \eqref{eq:scheme} on $[T_{j-1},T_j]$ with $\hZ_{u_i}=\psi^Z_j(\xi_i)$
  \EndFor
  \State form $e^{\sg}_j$ by \eqref{eq:sg-residual} and the batch loss $\frac1B\sum_b\sum_j\omega_j\abs{e^{\sg,(b)}_j}^2$
  \State update all parameters $\theta$ by one Adam step
\EndFor
\end{algorithmic}
\end{algorithm}

\subsection{Exercise rule and primal--dual bounds}\label{sec:certificates}

A trained network defines both an exercise rule and a dual martingale. Given a margin $\eex\ge0$, the rule exercises at the first date at which the payoff is positive and exceeds the predicted continuation value by at least $\eex$:
\[
\hat\tau=\min\bigl\{T_k,\ k<K:\ \Phi^{\mathrm f}_k>0,\ \Phi^{\mathrm f}_k\ge\hY_{k+1,T_k}+\eex\bigr\}\wedge T_K .
\]
The hedge heads define the martingale $\int_0^t\hZ_s^\top\dd W_s$, accumulated across all layers, and the lower and upper bounds are
\begin{equation}\label{eq:bounds}
L_0=\E\,\Phi^{\mathrm f}_{\hat\tau},\quad
U_0=\E\max_{1\le k\le K}\left(\Phi^{\mathrm f}_k-\int_0^{T_k}\hZ_s^\top\dd W_s\right),
\end{equation}
following the primal--dual construction of \cite{rogers2002,haugh2004}. Both bounds are estimated on independent test paths in a single forward pass. The asset position is recovered from $\hZ$ by $\hat\Delta_{u_i}=(\sigma(u_i,S_{u_i})^\top)^{-1}\hZ_{u_i}$ when $\sigma$ has full rank.
